\section{紧性、乘积与极小化}\label{sec:01-02}

\subsection{紧性：有限交、网与超滤子的统一}\label{subsec:1-2-1}

在欧氏空间里，Bolzano--Weierstrass 定理很容易使人把紧性理解为“每个序列都有收敛子列”。一般拓扑中，
序列只观察可数长的路线，而开覆盖可能同时包含不可数多个局部条件。下面的乘积首先暴露这个缺口。

\begin{example}[一个没有收敛子列的坐标序列]\label{ex:1-2-1-1}
令
\[
 X=\{0,1\}^{\mathcal P(\N)},\qquad x_n(A)=\1_A(n).
\]
任取子序列 \((x_{n_k})\)。令 \(A=\{n_{2k}:k\ge1\}\)，则
\[
 x_{n_k}(A)=
 \begin{cases}
  0,&k\text{ 为奇数},\\
  1,&k\text{ 为偶数}.
 \end{cases}
\]
因此该子序列在第 \(A\) 个坐标不收敛，也就不可能逐坐标收敛。到
\S\ref{subsec:1-2-2} 证明 Tychonoff 定理以后，我们会知道 \(X\) 在乘积拓扑下紧；在此之前只使用上面的
非收敛计算。
\end{example}

\begin{definition}[紧空间]\label{def:1-2-1-1}
拓扑空间 \(X\) 称为\emph{紧的}，若每个开覆盖都有有限子覆盖。闭集族具有\emph{有限交性质}，是指
其每个有限子族（包括空子族）的交非空；空子族的交约定为 \(X\)。
\end{definition}

\begin{proposition}[开覆盖与有限交性质]\label{proposition:1-2-1-cover-fip}
空间 \(X\) 紧，当且仅当每个具有有限交性质的闭集族都有非空交。这个约定包括空空间的情形。
\end{proposition}
\begin{proof}
设 \(X\) 紧，闭集族 \((F_i)_{i\in I}\) 的总交为空。则 \((X\setminus F_i)_{i\in I}\) 是开覆盖，
故有有限子覆盖。取补集可知相应有限多个 \(F_i\) 的交为空，所以该闭集族不具有有限交性质。

反过来，若开覆盖 \((U_i)_{i\in I}\) 没有有限子覆盖，则每个有限 \(E\subset I\) 都满足
\[
 \bigcap_{i\in E}(X\setminus U_i)
 =X\setminus\bigcup_{i\in E}U_i\ne\varnothing.
\]
于是闭集族 \((X\setminus U_i)\) 具有有限交性质。若假设中的闭集结论成立，其总交非空；但总交是
\(X\setminus\bigcup_iU_i=\varnothing\)，矛盾。
\end{proof}

有限交性质与滤子基不是同一个条件。前者只要求原集合族的每个有限交非空；后者还要求对任意两个基元素，
原集合族中已有另一个元素包含在二者之交内。例如在 \(X=\{1,2,3\}\) 上，
\[
  \mathcal C=\bigl\{\{1,2\},\{2,3\}\bigr\}
\]
具有有限交性质，却不是滤子基，因为 \(\mathcal C\) 中没有元素包含于 \(\{2\}\)。若 \(\mathcal C\)
具有有限交性质，则把它的所有有限交加入后，所得集合族便构成滤子基，并生成一个真滤子。这个加闭步骤正是
后面把闭集有限交性质转成滤子或超滤子时所用的操作。

\begin{definition}[相对紧]\label{def:1-2-1-relative-compact}
对子集 \(A\subset X\)，若其在环境空间 \(X\) 中的闭包 \(\overline A^{\,X}\) 紧，则称 \(A\) 在 \(X\)
中\emph{相对紧}。环境空间是定义的一部分；例如 \((0,1)\) 在 \(\R\) 中相对紧，而把 \((0,1)\) 自身作为
环境时，它的闭包仍为非紧空间 \((0,1)\)。
\end{definition}

\begin{proposition}[聚点的尾闭包表示]\label{proposition:1-2-1-tail-closure}
对网 \((x_\alpha)_{\alpha\in I}\)，令
\[
 K_\alpha=\overline{\{x_\beta:\beta\ge\alpha\}}.
\]
则该网的聚点集合为 \(\bigcap_{\alpha\in I}K_\alpha\)。
\end{proposition}
\begin{proof}
若 \(x\) 是聚点，给定 \(\alpha\) 与 \(x\) 的邻域 \(U\)，频繁进入性质给出
\(\beta\ge\alpha\) 使 \(x_\beta\in U\)。故 \(x\in K_\alpha\)，对每个 \(\alpha\) 都成立。

反过来，若 \(x\in\bigcap_\alpha K_\alpha\)，则对任意阈值 \(\alpha\) 与邻域 \(U\ni x\)，闭包定义给出
\(U\cap\{x_\beta:\beta\ge\alpha\}\ne\varnothing\)。所以网频繁进入每个 \(x\) 的邻域，\(x\) 是聚点。
\end{proof}

\begin{definition}[序列紧与可数紧]\label{def:1-2-1-3}
空间称为\emph{序列紧的}，若每个序列都有收敛子序列；称为\emph{可数紧的}，若每个可数开覆盖都有有限
子覆盖。一般拓扑空间中，这两个性质均不等同于紧性。
\end{definition}

\begin{theorem}[网紧性定理]\label{theorem:1-2-1-1}
空间 \(X\) 紧，当且仅当每个 \(X\) 中的网都有聚点；等价地，每个网都有收敛子网。
\end{theorem}
\begin{proof}
先设 \(X\) 紧。对网 \((x_\alpha)\) 考虑上一命题的尾闭包 \(K_\alpha\)。给定有限多个指标，取其共同后继
\(\gamma\)，则 \(K_\gamma\) 包含于相应有限交，故 \((K_\alpha)\) 具有有限交性质。紧性给出
\(x\in\bigcap_\alpha K_\alpha\)，上一命题说明 \(x\) 是聚点。由
\cref{theorem:1-1-3-1}，网有收敛子网。

反过来，设每个网都有聚点。令 \((F_i)_{i\in I}\) 是具有有限交性质的闭集族。若这个情形发生，则
\(X\ne\varnothing\)。令
\[
 J=\{(E,x):E\subset I\text{ 有限},\ x\in\bigcap_{i\in E}F_i\},
\]
并按有限集合坐标增大排序。有限交性质保证 \(J\ne\varnothing\)。两个指标的有限集合坐标取并，再从对应
非空交中取一点，得到共同后继，所以 \(J\) 有向。令网值 \(z_{(E,x)}=x\)，并取其聚点 \(z\)。

固定 \(i\in I\)。在任何有限集合坐标含 \(i\) 的指标之后，所有网值都落在 \(F_i\)。若 \(z\notin F_i\)，
则开邻域 \(X\setminus F_i\) 不可能被该尾部访问，与 \(z\) 是聚点矛盾。因此 \(z\in F_i\) 对每个 \(i\)
成立，即 \(\bigcap_iF_i\ne\varnothing\)。由有限交刻画，\(X\) 紧。这个反向构造使用全部可行对，不为每个
有限交预选一点，属于 \textsf{ZF} 证明。

“每个网有收敛子网”蕴含“每个网有聚点”由聚点--子网定理的逆向得到，故三种陈述等价。
\end{proof}

\begin{theorem}[超滤子紧性定理]\label{theorem:1-2-1-2}
在采用超滤子引理的语境中，\(X\) 紧，当且仅当 \(X\) 上每个超滤子至少收敛到一点。
\end{theorem}
\begin{proof}
先设 \(X\) 紧，\(\mathcal U\) 是超滤子。闭集族
\(\{\overline U:U\in\mathcal U\}\) 具有有限交性质，因为有限多个 \(U\) 的交仍属于真滤子且非空。
取
\[
 x\in\bigcap_{U\in\mathcal U}\overline U.
\]
对每个 \(V\in\mathcal N(x)\) 和 \(U\in\mathcal U\)，由 \(x\in\overline U\) 得 \(V\cap U\ne\varnothing\)，
所以 \(x\) 是 \(\mathcal U\) 的聚点。超滤子的聚点是极限，故 \(\mathcal U\to x\)。这一方向不需扩张
滤子，属于 \textsf{ZF} 证明。

反过来，设每个超滤子都有极限，令 \((F_i)\) 为具有有限交性质的闭集族。有限交组成滤子基；由超滤子引理，
其生成滤子可扩张为超滤子 \(\mathcal U\)。取 \(\mathcal U\to x\)。每个 \(F_i\in\mathcal U\)。若
\(x\notin F_i\)，则 \(X\setminus F_i\) 是 \(x\) 的邻域，也属于 \(\mathcal U\)，导致空集属于
\(\mathcal U\)。故 \(x\in\bigcap_iF_i\)，由有限交刻画可知 \(X\) 紧。这一方向使用 \textsf{UFL}。
\end{proof}

\begin{proposition}[紧集在 Hausdorff 空间中闭]\label{proposition:1-2-1-3}
若 \(K\subset X\) 紧且 \(X\) Hausdorff，则 \(K\) 闭。Hausdorff 不能只替换为 \(T_1\)。
\end{proposition}
\begin{proof}
取 \(x\in X\setminus K\)。令 \(\mathscr U\) 为所有满足下列条件的开集 \(U\)：存在 \(x\) 的开邻域
\(V\) 使 \(U\cap V=\varnothing\)。Hausdorff 性保证每个 \(y\in K\) 属于某个 \(U\in\mathscr U\)，
所以 \(\mathscr U\) 覆盖 \(K\)。紧性给出有限子覆盖 \(U_1,\ldots,U_m\)。对这有限多个 \(U_j\)，分别取
\(x\) 的开邻域 \(V_j\) 使 \(U_j\cap V_j=\varnothing\)。则 \(V=\bigcap_jV_j\) 是 \(x\) 的开邻域，
并与 \(\bigcup_jU_j\supset K\) 不交。因此 \(X\setminus K\) 开，\(K\) 闭。

反例取无限集合上的余有限拓扑。它是 \(T_1\) 且紧：任取开覆盖，先取一个非空成员，其补集有限，再为剩余
有限多个点各取一个覆盖成员即可。任一无限真子集在子空间余有限拓扑下也紧，却不是环境空间中的闭集。
\end{proof}

\begin{example}[递减闭集网]\label{ex:1-2-1-2}
若 \(X\) 紧，\((F_\alpha)\) 是按指标递减的非空闭集网，则这些闭集具有有限交性质，所以
\(\bigcap_\alpha F_\alpha\ne\varnothing\)。可算模型是
\(\bigcap_n[0,1/n]=\{0\}\)。若删去紧性，\(\R\) 中 \(F_n=[n,\infty)\) 每个都非空闭且递减，总交却为空：
约束的满足点逃向无穷。
\end{example}

\begin{theorem}[度量空间中的序列判别]\label{theorem:1-2-1-4}
度量空间 \(X\) 紧，当且仅当 \(X\) 序列紧。
\end{theorem}
\begin{proof}
若 \(X=\varnothing\)，则开覆盖紧性与序列紧性都成立。以下设 \(X\ne\varnothing\)。
设 \(X\) 紧，\((x_n)\) 是序列。把它看成网，网紧性定理给出聚点 \(x\)。递归取严格递增整数 \(n_k\)，使
\[
 n_k>n_{k-1},\qquad d(x_{n_k},x)<1/k.
\]
聚点性质保证每一步都有可选的 \(n_k\)。于是 \(x_{n_k}\to x\)。这份递归抽取使用 \textsf{DC}。

反过来设 \(X\) 序列紧。先证全有界。若存在 \(\varepsilon>0\) 使 \(X\) 不能由有限多个
\(\varepsilon\)-球覆盖，可递归选择 \(x_{n+1}\) 不属于
\(\bigcup_{j\le n}B(x_j,\varepsilon)\)。所得序列两两距离至少 \(\varepsilon\)，所以没有 Cauchy 子列，
更没有收敛子列，矛盾。

对每个 \(m\ge1\)，取有限 \(1/m\)-网 \(D_m\)，则
\[
 \mathcal B=\{B(a,r):a\in\bigcup_mD_m,\ r\in\Q,\ r>0\}
\]
是可数基。事实上，若 \(x\in U\) 且 \(U\) 开，取 \(\rho>0\) 使 \(B(x,\rho)\subset U\)，再取
\(m\) 使 \(2/m<\rho\) 及 \(a\in D_m\) 使 \(d(a,x)<1/m\)，并选有理
\(r\) 满足 \(1/m<r<\rho-1/m\)。于是 \(x\in B(a,r)\subset U\)。

给定开覆盖 \(\mathscr U\)，对每个被某个覆盖成员包含的基元素，选取一个这样的覆盖成员。所有选出的成员
构成可数子覆盖；把它编号为 \(U_1,U_2,\ldots\)。若它没有有限子覆盖，逐 \(n\) 取
\[
 x_n\in X\setminus\bigcup_{j=1}^nU_j.
\]
取收敛子列 \(x_{n_k}\to x\)，并选 \(m\) 使 \(x\in U_m\)。当 \(k\) 充分大时
\(x_{n_k}\in U_m\)，而 \(n_k\ge m\) 时定义又给 \(x_{n_k}\notin U_m\)，矛盾。因此原开覆盖有有限
子覆盖。逐尺度取有限网、从可数子覆盖取反例点和递归抽取所使用的原则不超过通常的
\textsf{CC}/\textsf{DC}；这里不主张这是命题的最优集合论强度。
\end{proof}

\begin{example}[紧性、序列紧性与全有界性]\label{ex:1-2-1-3}
在度量空间中，上一证明说明开覆盖紧性和序列紧性重新一致。结合完备性与全有界性还可得到第三种等价刻画；
其证明留作本小节习题，以便把一般拓扑中的区别与度量特例分开。

全有界性不能只换成有界性。例如 \(\ell^2\) 的闭单位球完备且有界，但标准正交向量满足
\(\norm{e_m-e_n}=\sqrt2\)（\(m\ne n\)）。因而半径小于 \(\sqrt2/2\) 的有限多个球不能覆盖该闭单位球；
它不是全有界的，也就不是紧的。
\end{example}

\begin{figure}[htbp]
\centering
\begin{tikzpicture}[node distance=1.35cm and 1.7cm,>=Stealth,every node/.style={align=center}]
 \node[book strong node] (a) {紧性\\有限子覆盖};
 \node[book node,left=of a] (b) {闭集 FIP\\交非空};
 \node[book node,right=of a] (c) {网有聚点\\有收敛子网};
 \node[book node,below=of a] (d) {每个超滤子\\有极限};
 \draw[book arrow,<->] (a)--node[above]{补集}(b);
 \draw[book arrow,<->] (a)--node[above]{\textsf{ZF}}(c);
 \draw[book arrow] (a)--node[left]{\textsf{ZF}}(d);
 \draw[book dashed arrow] (d)--node[right]{反向用 \textsf{UFL}}(a);
\end{tikzpicture}
\caption{紧性的四种语言以及超滤子反向证明所需的扩张原则。}
\label{fig:1-2-1-compactness-equivalences}
\end{figure}

紧性给出超滤子极限的存在，却不保证唯一；唯一性仍由 Hausdorff 性提供。相对紧也始终相对于环境空间定义。
在一般拓扑中，序列紧与可数紧都不能替代紧性；度量结构提供的可数基才使序列重新足够。

\begin{exercise}
证明度量空间紧当且仅当它完备且全有界。两个方向都要写出，并说明从全有界序列抽取 Cauchy 子列时
怎样逐层选择有限网中的嵌套无限子列。
\end{exercise}
\begin{exercise}
不使用 Tychonoff 定理，证明有限个紧空间的乘积紧。
\end{exercise}
\begin{exercise}
证明紧空间到 Hausdorff 空间的连续双射是同胚。
\end{exercise}
\begin{exercise}
令 \(\omega_1\) 为第一不可数序数，并给序数区间 \([0,\omega_1)\) 配置序拓扑。证明每个可数子集
都有严格小于 \(\omega_1\) 的上界，并据此证明该空间既可数紧又序列紧，但不紧。最后证明它是
Hausdorff 空间，并指出论证中何处使用了“可数个可数序数的上确界仍是可数序数”。
本题是集合论拓扑的选读旁例；序数、\(\omega_1\)、序拓扑以及引号内的集合论结论在本题中作为外部
定义与待证引理，本章正文后续不依赖本题。
\end{exercise}
\begin{exercise}
证明在紧 Hausdorff 空间中，网收敛到 \(x\) 当且仅当 \(x\) 是它的唯一聚点。指出“唯一聚点蕴含收敛”
使用紧性，而排除第二聚点使用 Hausdorff 性。
\end{exercise}

这些等价刻画处理的是单个空间。函数族与坐标空间把许多局部紧性问题同时放在各个坐标上，因而接下来要
判断：逐坐标的有限观测能否重新拼成整体紧性。

\subsection{乘积拓扑、逐点收敛与 Tychonoff 定理}\label{subsec:1-2-2}

一个函数可以看作其全部函数值组成的坐标族。若一次实验只能检查有限多个输入点，那么合适的邻域也只应
限制有限多个坐标。允许一次限制所有坐标会得到更细的盒拓扑；下面的失败例将说明，那种拓扑既不描述点态
收敛，也不保留乘积紧性。

\begin{remark}
有限乘积紧性可以反复使用两因子情形得到；任意乘积则把选择原则带入拓扑核心。超滤子证明会把两件事分开：
因子紧性产生坐标极限，而有限柱只负责把坐标收敛拼成乘积收敛。
\end{remark}

\begin{definition}[乘积拓扑]\label{def:1-2-2-1}
对空间族 \((X_i)_{i\in I}\)，乘积 \(X=\prod_{i\in I}X_i\) 上的\emph{乘积拓扑}是使每个坐标投影
\(\pi_i:X\to X_i\) 连续的最弱拓扑。它的基本开集为
\[
 \bigcap_{i\in F}\pi_i^{-1}(U_i),
\]
其中 \(F\subset I\) 有限，且每个 \(U_i\subset X_i\) 开。
\end{definition}

\begin{definition}[逐点收敛拓扑]\label{def:1-2-2-2}
函数空间 \(Y^S=\prod_{s\in S}Y\) 上的乘积拓扑称为\emph{逐点收敛拓扑}。
\end{definition}

\begin{definition}[柱集]\label{def:1-2-2-3}
只由有限多个坐标条件决定的集合称为\emph{柱集}。坐标条件均开时，所得开柱构成乘积拓扑的基。
\end{definition}

\begin{theorem}[逐坐标收敛判别]\label{theorem:1-2-2-1}
网 \(x_\alpha\in\prod_iX_i\) 收敛到 \(x\)，当且仅当对每个 \(i\in I\)，
\(\pi_i(x_\alpha)\to\pi_i(x)\)。
\end{theorem}
\begin{proof}
若 \(x_\alpha\to x\)，坐标投影连续，所以每个坐标网收敛。

反过来，设每个坐标网收敛。取 \(x\) 的基本邻域
\[
 W=\bigcap_{i\in F}\pi_i^{-1}(U_i),
\]
其中 \(F\) 有限且 \(U_i\) 是 \(x_i\) 的邻域。对每个 \(i\in F\)，存在阈值 \(\alpha_i\)，使
\(\alpha\ge\alpha_i\) 时 \(\pi_i(x_\alpha)\in U_i\)。有限次使用有向性得到共同后继 \(\alpha_0\)。
当 \(\alpha\ge\alpha_0\) 时所有有限坐标条件同时成立，故 \(x_\alpha\in W\)。每个邻域包含一个这样的
基本邻域，所以 \(x_\alpha\to x\)。
\end{proof}

\begin{proposition}[滤子的逐坐标收敛判别]\label{proposition:1-2-2-filter-coordinate}
乘积上的滤子 \(\mathcal F\) 收敛到 \(x=(x_i)\)，当且仅当
\((\pi_i)_*\mathcal F\to x_i\) 对每个 \(i\) 成立。
\end{proposition}
\begin{proof}
正向由坐标投影的连续性和滤子连续性判别得到。反过来，取 \(x\) 的基本柱邻域
\(W=\bigcap_{i\in F}\pi_i^{-1}(U_i)\)。坐标像滤子收敛给出
\(\pi_i^{-1}(U_i)\in\mathcal F\)。滤子对有限交封闭，故 \(W\in\mathcal F\)。因此每个 \(x\) 的邻域
属于 \(\mathcal F\)，即 \(\mathcal F\to x\)。
\end{proof}

\begin{example}[函数族的点态极限]\label{ex:1-2-2-1}
若每个 \(s\in S\) 的函数值都落在紧集 \(K_s\)，函数族 \(\mathcal A\) 可视为
\(\prod_sK_s\) 的子集。Tychonoff 定理将保证环境乘积紧，所以 \(\mathcal A\) 中任意网有逐点收敛子网。
极限先落在 \(\prod_sK_s\)；要使它仍属于 \(\mathcal A\)，还需 \(\mathcal A\) 对逐点极限封闭。
要把逐点收敛升级为一致收敛，则需要 \S\ref{subsec:1-3-2} 的等度连续性。

例如令 \(S=[0,1]\)、\(K_s=[0,1]\)，并令 \(\mathcal M\) 为所有单调不减映射
\(f:[0,1]\to[0,1]\) 的集合。若网 \(f_\alpha\in\mathcal M\) 逐点收敛到 \(f\)，则对
\(x\le y\) 有
\[
  f(x)=\lim_\alpha f_\alpha(x)\le \lim_\alpha f_\alpha(y)=f(y).
\]
所以 \(\mathcal M\) 在乘积中闭，从而在逐点拓扑下紧。这个结论并不声称极限连续；例如连续函数
\(f_n(x)=\min\{nx,1\}\) 逐点收敛到在 \(0\) 处不连续的单调函数。
\end{example}

\begin{example}[可数乘积度量的构造]\label{ex:1-2-2-2}
若 \(X_n\) 有度量 \(d_n\le1\)，则
\[
 d(x,y)=\sum_{n=1}^{\infty}2^{-n}d_n(x_n,y_n)
\]
定义一个度量：级数由 \(\sum_n2^{-n}=1\) 控制；若 \(d(x,y)=0\)，则每个非负项
\(2^{-n}d_n(x_n,y_n)\) 都为零，故 \(x_n=y_n\) 对所有 \(n\) 成立；三角不等式可对各坐标的
三角不等式逐项求和。\cref{proposition:1-2-2-3} 将进一步证明它诱导乘积拓扑。
\end{example}

\begin{example}[盒拓扑的双重失败]\label{ex:1-2-2-3}
若基本开集可以限制所有坐标，得到盒拓扑。在 \(\R^\N\) 中令
\(x^{(k)}_n=1/k\)。对每个固定 \(n\)，有 \(x^{(k)}_n\to0\)，所以它在乘积拓扑中趋于零；但它从不进入
盒邻域
\[
 \prod_{n=1}^{\infty}(-1/n,1/n),
\]
因为对任意 \(k\)，取 \(n>k\) 即有 \(1/k\notin(-1/n,1/n)\)。

盒拓扑也破坏紧性。考虑 \(D=\{(t,t,\ldots):t\in[0,1]\}\subset[0,1]^\N\)。若 \(x\notin D\)，
存在两个坐标 \(x_m\ne x_n\)。令 \(\eta=|x_m-x_n|/3\)，并在第 \(m,n\) 个坐标分别限制到
\[
  (x_m-\eta,x_m+\eta)\cap[0,1],\qquad
  (x_n-\eta,x_n+\eta)\cap[0,1].
\]
两区间不交；其余坐标取 \([0,1]\)，所得盒邻域不可能含对角点，故 \(D\) 闭。
对固定 \(t\)，在第 \(n\) 个坐标取相对邻域
\((t-1/n,t+1/n)\cap[0,1]\)。若对角点 \((s,s,\ldots)\) 属于这些区间的盒积，则
\[
  |s-t|<\frac1n\quad(n\ge1),
\]
从而 \(s=t\)。故该盒积与 \(D\) 只交于 \((t,t,\ldots)\)，所以 \(D\) 在子空间拓扑下离散。
单点族 \(\{\{p\}:p\in D\}\) 是 \(D\) 的开覆盖且没有有限子覆盖，故 \(D\) 不紧。若盒乘积紧，
其闭子集 \(D\) 应紧，矛盾。因此盒乘积不紧。
\end{example}

\begin{remark}[不可数乘积中的 Borel--柱差异]\label{rem:1-2-2-borel-cylinder}
为比较不可数乘积上的两种生成过程，在此使用以下集合概念：\(\sigma\)-代数是对补集与可数并封闭的
集合族，Borel \(\sigma\)-代数由全部开集生成。\S\ref{subsec:1-4-2} 将系统讨论这些定义。

令 \(I\) 不可数，\(X=\{0,1\}^{I}\)。称 \(E\subset X\) 依赖坐标集 \(J\subset I\)，若任意两点在
\(J\) 上坐标相同便同时属于或同时不属于 \(E\)。所有依赖某个可数坐标集的集合组成 \(\sigma\)-代数：
补集保留同一坐标集；可数并依赖对应可数坐标集之并。它包含全部有限柱，所以有限柱生成的乘积
\(\sigma\)-代数中，每个集合至多依赖可数多个坐标。

然而开集
\[
 X\setminus\{0^I\}=\bigcup_{i\in I}\{x:x_i=1\}
\]
不依赖任何可数 \(J\subset I\)：取 \(i\in I\setminus J\)，全零点与只在第 \(i\) 坐标为 \(1\) 的点在
\(J\) 上相同，却分别不属于和属于该开集。因此乘积拓扑的 Borel \(\sigma\)-代数可以严格大于有限柱生成的
\(\sigma\)-代数。这里使用“可数个可数集之并可数”时处在通常的 \textsf{CC}/ZFC 语境中。
\end{remark}

\begin{theorem}[Tychonoff 定理]\label{theorem:1-2-2-2}
若每个 \(X_i\) 紧，则 \(\prod_iX_i\) 在乘积拓扑下紧。
\end{theorem}
\begin{proof}
若某个因子为空，乘积为空，结论成立。以下设每个因子非空。

先在通常的 ZFC 语境中证明一般版本。给定乘积 \(X=\prod_iX_i\) 上的超滤子 \(\mathcal U\)，
\cref{proposition:1-1-3-4} 说明 \((\pi_i)_*\mathcal U\) 是 \(X_i\) 上超滤子。因 \(X_i\) 紧，
\cref{theorem:1-2-1-2} 给出非空极限集
\[
 L_i=\{x_i\in X_i:(\pi_i)_*\mathcal U\to x_i\}.
\]
对所有 \(i\) 同时选取 \(x_i\in L_i\)，得到 \(x=(x_i)\in X\)。由滤子的逐坐标判别，
\(\mathcal U\to x\)。所以 \(X\) 上每个超滤子都有极限，超滤子紧性定理给出 \(X\) 紧。这份证明既使用
超滤子扩张，也同时选择坐标极限，按 \textsf{AC} 记录。

若每个 \(X_i\) 还 Hausdorff，则每个 \(L_i\) 至多含一点；紧性又保证它非空，所以坐标极限由
\((\pi_i)_*\mathcal U\) 唯一确定，不需从一族多元素集合中同时选代表。在这种情形，采用
\textsf{UFL} 足以完成上述证明。本书后面使用的 \([0,1]^S\) 和紧值函数盒均属于紧 Hausdorff 情形。
\end{proof}

由 Tychonoff 定理，\cref{ex:1-2-1-1} 中的
\(\{0,1\}^{\mathcal P(\N)}\) 确实紧；结合该例中的坐标计算，它是一个紧而非序列紧的空间。

\begin{proposition}[可数乘积的第二可数性与可度量化]\label{proposition:1-2-2-3}
可数个第二可数空间的乘积第二可数；可数个可度量空间的乘积可度量。若每个因子紧且可度量，则乘积紧且
可度量，并可用对角子序列证明序列紧性。
\end{proposition}
\begin{proof}
对每个 \(n\)，取可数基 \(\mathcal B_n\)，并把 \(X_n\) 添入该基。所有只在某个有限坐标集
\(F\subset\N\) 上取 \(\mathcal B_n\) 成员、其余坐标不限制的柱集构成可数族；每个基本柱都含有
其中一个成员，所以乘积第二可数。

若 \(d_n\) 是 \(X_n\) 的度量，令 \(\rho_n=\min\{1,d_n\}\) 并定义
\[
 d(x,y)=\sum_{n=1}^{\infty}2^{-n}\rho_n(x_n,y_n).
\]
半径小于 \(1\) 的 \(\rho_n\)-球与同半径的 \(d_n\)-球相同，所以 \(\rho_n\) 与 \(d_n\) 诱导同一拓扑。
而且
\[
 \rho_n(x,z)\le \min\{1,d_n(x,y)+d_n(y,z)\}
 \le \rho_n(x,y)+\rho_n(y,z),
\]
所以 \(\rho_n\) 是度量，级数逐项给出 \(d\) 的三角不等式。若 \(d(x,y)=0\)，每项均为零，故 \(x=y\)。
固定 \(x\) 并给定 \(\varepsilon>0\)。选 \(N\) 使
\(\sum_{n>N}2^{-n}<\varepsilon/2\)，再令
\[
 W=\bigcap_{n=1}^N
   \pi_n^{-1}\!\left(B_{\rho_n}(x_n,\varepsilon/2)\right).
\]
若 \(y\in W\)，则前段和严格小于
\((\varepsilon/2)\sum_{n=1}^N2^{-n}<\varepsilon/2\)，尾段和也小于
\(\varepsilon/2\)，故 \(W\subset B_d(x,\varepsilon)\)。
反过来，设 \(x\) 的基本柱邻域只限制有限坐标集 \(F\subset\N\)。对每个 \(n\in F\)，取
\(r_n>0\) 使 \(B_{\rho_n}(x_n,r_n)\) 包含于相应坐标邻域。若 \(F\ne\varnothing\)，取
\[
 0<\delta<\min_{n\in F}2^{-n}\min\{1,r_n\}.
\]
若 \(d(x,y)<\delta\)，则对每个 \(n\in F\)，
\(2^{-n}\rho_n(x_n,y_n)<\delta\)，所以 \(y_n\) 落入相应坐标球。故 \(d\) 诱导乘积拓扑。
当 \(F=\varnothing\) 时柱邻域就是全空间，反向包含显然成立。

若每个因子紧可度量，Tychonoff 给出乘积紧。也可直接对任意序列逐坐标抽取：先取第一坐标收敛的子序列，
再从中取第二坐标收敛的子序列，如此递归；取第 \(k\) 层子序列的第 \(k\) 项作为对角序列。每个固定坐标
最终位于相应层子序列中，故逐坐标收敛，再由乘积度量收敛。递归抽取使用 \textsf{DC}/\textsf{CC}。
\end{proof}

\begin{proposition}[映入乘积的连续性]\label{proposition:1-2-2-4}
映射 \(F:Z\to\prod_iX_i\) 连续，当且仅当每个坐标映射 \(\pi_i\circ F\) 连续。
\end{proposition}
\begin{proof}
正向由连续映射复合得到。反过来，基本柱
\(W=\bigcap_{i\in E}\pi_i^{-1}(U_i)\) 的逆像为
\[
 F^{-1}(W)=\bigcap_{i\in E}(\pi_i\circ F)^{-1}(U_i),
\]
这是有限多个开集之交，因而开。基本柱的任意并给出全部开集，所以 \(F\) 连续。
\end{proof}

\begin{figure}[htbp]
\centering
\begin{tikzpicture}[node distance=1.35cm and 1.45cm,>=Stealth,every node/.style={align=center}]
 \node[book strong node] (u) {乘积超滤子 \(\mathcal U\)};
 \node[book node,right=of u] (p) {坐标超滤子\\\((\pi_i)_*\mathcal U\)};
 \node[book node,right=of p] (l) {坐标极限\\\(x_i\in L_i\)};
 \node[book node,below=of l] (x) {乘积点 \(x=(x_i)\)};
 \node[book node,left=of x] (c) {有限柱检验};
 \draw[book arrow] (u)--(p);
 \draw[book arrow] (p)--node[above]{因子紧}(l);
 \draw[book arrow] (l)--node[right]{一般版需选择}(x);
 \draw[book arrow] (u)|-(c);
 \draw[book arrow] (c)--node[below]{有限交}(x);
\end{tikzpicture}
\caption{坐标极限的存在与有限柱的收敛检验是两个不同步骤。}
\label{fig:1-2-2-tychonoff-coordinate-choice}
\end{figure}

乘积拓扑只给逐坐标控制；连续性和一致收敛需要额外结构。一般紧空间版、紧 Hausdorff 版和可数紧可度量版
的证明采用不同原则，不能把其中一份证明无条件搬到另一种因子类别。

\begin{exercise}
证明任意乘积中的坐标投影是开映射。
\end{exercise}
\begin{exercise}
在乘积非空的前提下，证明乘积 Hausdorff 当且仅当每个因子 Hausdorff；再用某个因子为空的情形说明为何
需要“乘积非空”。
\end{exercise}
\begin{exercise}
若每个 \((X_n,d_n)\) 完备，证明上述乘积度量完备。
\end{exercise}
\begin{exercise}
完整证明不可数 \(I\) 时，\(\{0,1\}^{I}\) 的有限柱生成 \(\sigma\)-代数严格小于乘积拓扑的 Borel
\(\sigma\)-代数。
\end{exercise}

乘积紧性负责产生候选极限，却不保证某个数值泛函在极限处保留正确的不等式。变分问题因此还需要一种
只控制“向下跳跃”的极限条件。

\subsection{半连续性、\texorpdfstring{\(\liminf\)}{liminf} 与紧性原理}
\label{subsec:1-2-3}

紧性能够阻止近似极小点逃走，却不能保证极限仍然极小。反过来，正确方向的极限不等式能够识别一个已经
存在的候选点，却不能自行产生候选。先看两种彼此独立的失败。

\begin{remark}
Weierstrass 极值定理通常为连续函数陈述，但直接法只需要禁止函数值在极限处向下跳，这正是下半连续。
把“候选点存在”与“候选点保持极小不等式”分开以后，紧性与半连续性的职责便不再混在一起。
\end{remark}

\begin{example}[两种最小化失败]\label{ex:1-2-3-3}
连续函数 \(f(x)=e^x\) 在 \(\R\) 上的下确界为 \(0\)，但没有最小点；点列趋向 \(-\infty\) 时函数值
趋零，候选点逃出了空间。第二种失败发生在紧集上：令
\[
 g(0)=1,\qquad g(x)=x\quad(0<x\le1).
\]
则 \(\inf_{[0,1]}g=0\)，仍没有最小点。这里空间不允许逃逸，失败来自
\(g(0)>\liminf_{x\downarrow0}g(x)\)。紧性与正确的极限不等式必须同时出现。
\end{example}

\begin{definition}[紧扩展实直线]\label{def:1-2-3-extended-line}
令 \(\overline\R=[-\infty,+\infty]\) 带次序拓扑。映射
\[
 h(t)=\frac{2}{\pi}\arctan t\quad(t\in\R),\qquad
 h(-\infty)=-1,\quad h(+\infty)=1
\]
是 \(\overline\R\) 到 \([-1,1]\) 的同胚。因此 \(\overline\R\) 是紧 Hausdorff 空间。
\end{definition}
\begin{proof}
\(h\) 在 \(\R\) 上严格递增、连续，并把 \(\R\) 双射到 \((-1,1)\)。当
\(t\to\pm\infty\) 时 \(h(t)\to\pm1\)，所以端点延拓连续。其逆在 \((-1,1)\) 上为
\(\tan(\pi s/2)\)，并在 \(s\to\pm1\) 时趋于相应无穷端点，故逆映射也连续。
\end{proof}

\begin{definition}[下半连续]\label{def:1-2-3-1}
函数 \(f:X\to\overline\R\) 称为\emph{下半连续}，若对每个 \(a\in\R\)，集合
\(\{x:f(x)>a\}\) 开。等价地，每个次水平集 \(\{x:f(x)\le a\}\) 闭。上半连续用相反不等式定义。
\end{definition}

\begin{proposition}[上图判别]\label{proposition:1-2-3-epigraph}
对 \(f:X\to\overline\R\)，令
\[
 \epi(f)=\{(x,t)\in X\times\R:f(x)\le t\}.
\]
则 \(f\) 下半连续，当且仅当 \(\epi(f)\) 在 \(X\times\R\) 中闭。
\end{proposition}
\begin{proof}
若上图闭，固定 \(a\in\R\)，连续映射 \(i_a:X\to X\times\R\)、\(i_a(x)=(x,a)\) 满足
\[
 i_a^{-1}(\epi(f))=\{x:f(x)\le a\},
\]
所以每个次水平集闭，\(f\) 下半连续。

反过来设 \(f\) 下半连续，取 \((x,t)\notin\epi(f)\)，即 \(f(x)>t\)。选实数 \(a\) 满足
\(t<a<f(x)\)。集合
\[
 \{y:f(y)>a\}\times(-\infty,a)
\]
是 \((x,t)\) 的开邻域；其中每个 \((y,s)\) 满足 \(f(y)>a>s\)，故不属于上图。因此上图的补集开，
上图闭。
\end{proof}

\begin{example}[边界值决定上图是否闭]\label{ex:1-2-3-step-epigraph}
在 \(\R\) 上，\(f=\1_{(0,\infty)}\) 下半连续：当 \(0\le a<1\) 时
\(\{f>a\}=(0,\infty)\)，其他阈值的超水平集是 \(\R\) 或空集。改令
\(g=\1_{[0,\infty)}\)，则 \(x_n=-1/n\to0\)，而
\[
 g(0)=1>\liminf_ng(x_n)=0.
\]
等价地，\((x_n,0)\in\epi(g)\)，其极限 \((0,0)\) 不在 \(\epi(g)\)。这个计算固定了下半连续不等式的
方向。
\end{example}

\begin{definition}[网的 \(\liminf\) 与 \(\limsup\)]\label{def:1-2-3-2}
对扩展实数网 \((a_\alpha)\)，定义
\[
 \liminf_\alpha a_\alpha
 =\sup_{\alpha_0}\inf_{\alpha\ge\alpha_0}a_\alpha,\qquad
 \limsup_\alpha a_\alpha
 =\inf_{\alpha_0}\sup_{\alpha\ge\alpha_0}a_\alpha.
\]
\end{definition}

\begin{proposition}[尾部聚点与 \(\liminf\)]\label{proposition:1-2-3-tail-cluster}
把 \((a_\alpha)\) 看成紧空间 \(\overline\R\) 中的网，并令
\[
 K_{\alpha_0}=\overline{\{a_\alpha:\alpha\ge\alpha_0\}},
 \qquad C=\bigcap_{\alpha_0}K_{\alpha_0}.
\]
则 \(C\) 是非空紧集，恰为该网的聚点集，而且
\[
 \min C=\liminf_\alpha a_\alpha,\qquad
 \max C=\limsup_\alpha a_\alpha.
\]
\end{proposition}
\begin{proof}
每个尾闭包都非空且闭。给定有限多个指标 \(\alpha_1,\ldots,\alpha_m\)，取共同后继
\(\gamma\)，则
\(K_\gamma\subset\bigcap_{j=1}^mK_{\alpha_j}\)；因此这些尾闭包具有有限交性质。扩展实直线紧，故
\(C\ne\varnothing\) 且紧。由
\cref{proposition:1-2-1-tail-closure}，\(C\) 正是聚点集。紧的非空扩展实数子集有最小值和最大值。
具体地，若非空闭集 \(K\subset\overline\R\) 且 \(b=\inf K\)，则 \(b=+\infty\) 时
\(K=\{+\infty\}\)；\(b=-\infty\) 时每个
\([-\infty,t)\) 都与 \(K\) 相交；\(b\in\R\) 时每个
\((b-\varepsilon,b+\varepsilon)\) 都与 \(K\) 相交。因此在所有情形都有 \(b\in\overline K=K\)；
上确界同理。

记 \(b_{\alpha_0}=\inf_{\alpha\ge\alpha_0}a_\alpha\) 和
\(\ell=\sup_{\alpha_0}b_{\alpha_0}\)。若
\(S_{\alpha_0}=\{a_\alpha:\alpha\ge\alpha_0\}\)，则
\(S_{\alpha_0}\subset[b_{\alpha_0},+\infty]\)。右侧闭，故
\(K_{\alpha_0}=\overline{S_{\alpha_0}}\subset[b_{\alpha_0},+\infty]\)，而
\(S_{\alpha_0}\subset K_{\alpha_0}\) 又给出反向的下确界不等式。因此闭包不改变下确界；上段论证还说明
该下确界属于闭包，故
\[
  \min K_{\alpha_0}=\min\overline{S_{\alpha_0}}
  =\inf S_{\alpha_0}=b_{\alpha_0}.
\]
所以 \(C\subset K_{\alpha_0}\) 给出 \(\min C\ge b_{\alpha_0}\)，进而
\(\min C\ge\ell\)。

若 \(\min C>\ell\)，取 \(t\) 使 \(\ell<t<\min C\)。由于
\(b_{\alpha_0}\le\ell<t\)，每个尾闭包都与闭区间 \([-\infty,t]\) 相交。更具体地，给定
有限多个指标 \(\alpha_1,\ldots,\alpha_m\)，取共同后继 \(\gamma\)；则
\(K_\gamma\subset\bigcap_jK_{\alpha_j}\)，而
\(\min K_\gamma=b_\gamma\le\ell<t\)，所以
\(K_\gamma\cap[-\infty,t]\ne\varnothing\)。这就证明闭集族
\[
 \{K_{\alpha_0}\}_{\alpha_0}\cup\{[-\infty,t]\}
\]
具有有限交性质；紧性给出 \(C\cap[-\infty,t]\ne\varnothing\)，与 \(\min C>t\) 矛盾。因此
\(\min C=\ell=\liminf_\alpha a_\alpha\)。最大值的论证把次序反向即可。
\end{proof}

\begin{definition}[次水平紧与等强制]\label{def:1-2-3-3}
函数 \(F:X\to\overline\R\) 称为\emph{次水平紧的}（inf-compact），若每个
\(\{F\le c\}\)、\(c\in\R\)，都紧。函数族 \((F_\lambda)\) 称为\emph{等强制的}，若对每个 \(c\in\R\)，
\[
 \bigcup_\lambda\{F_\lambda\le c\}
\]
相对紧。前者描述单个函数，后者统一控制一族函数。
\end{definition}

\begin{example}[可计算的等强制族]\label{ex:1-2-3-equicoercive}
在 \(\R\) 上令 \(F_n(x)=x^2+1/n\)。当 \(c<0\) 时所有次水平集均为空；当 \(c\ge0\) 时
\[
 \bigcup_n\{F_n\le c\}\subset[-\sqrt c,\sqrt c],
\]
所以 \((F_n)\) 等强制。若 \(\sup_nF_n(x_n)<\infty\)，则 \((x_n)\) 有界，因而相对紧。这个结论只产生
候选点；要比较变化中的 \(F_n\) 与某个极限泛函的值，还需额外的 \(\liminf\) 条件。相应的
\(\Gamma\)-收敛理论不在本节使用。
\end{example}

\begin{example}[开集指标与闭约束]\label{ex:1-2-3-1}
若 \(U\) 开，则 \(\1_U\) 下半连续；若 \(C\) 闭，则 \(\1_C\) 上半连续。对变分问题，更有用的闭约束
能量是
\[
 I_C(x)=
 \begin{cases}
  0,&x\in C,\\
  +\infty,&x\notin C.
 \end{cases}
\]
它的有限次水平集在 \(a\ge0\) 时为 \(C\)，在 \(a<0\) 时为空，所以 \(I_C\) 下半连续当且仅当 \(C\) 闭。
\end{example}

\begin{example}[有限维模型中的弱下半连续性]\label{ex:1-2-3-2}
先把本例所用术语就地定义：\(\R^d\) 的弱拓扑
\(\sigma(\R^d,(\R^d)^*)\) 是使每个线性泛函 \(\ell:\R^d\to\R\) 连续的最弱拓扑。记该拓扑为
\(\tau_w\)，欧氏拓扑为 \(\tau_e\)。每个线性泛函在 \(\tau_e\) 下连续，故
\(\tau_w\subset\tau_e\)。反过来，标准坐标泛函 \(p_j(x)=x_j\) 属于 \((\R^d)^*\)，而有限交
\[
 \bigcap_{j=1}^d p_j^{-1}\bigl((x_j-\varepsilon,x_j+\varepsilon)\bigr)
\]
构成欧氏拓扑的一套邻域基，故 \(\tau_e\subset\tau_w\)。于是 \(\tau_w=\tau_e\)。
Cauchy--Schwarz 不等式给出
\[
  \norm{x}=\sup_{\norm{u}\le1}\langle x,u\rangle,\qquad
  \{x:\norm{x}\le r\}
  =\bigcap_{\norm{u}\le1}\{x:\langle x,u\rangle\le r\},\qquad r\ge0.
\]
Cauchy--Schwarz 给出
\(\sup_{\norm{u}\le1}\langle x,u\rangle\le\norm{x}\)；当 \(x\ne0\) 时取
\(u=x/\norm{x}\) 即取到反向不等式，而 \(x=0\) 时等号显然。第二式于是由
\(\norm{x}\le r\) 当且仅当所有单位球内的 \(u\) 都满足
\(\langle x,u\rangle\le r\) 得到。
右端是闭半空间之交，所以闭球是弱闭集，范数因而弱下半连续。这个有限维论证只使用欧氏内积。
一般赋范空间中的相应闭半空间表示需要 Hahn--Banach 分离，见第~3~章。
\end{example}

\begin{theorem}[半连续性的网判别]\label{theorem:1-2-3-1}
函数 \(f:X\to\overline\R\) 下半连续，当且仅当对每个 \(x_\alpha\to x\)，
\[
 f(x)\le\liminf_\alpha f(x_\alpha).
\]
\end{theorem}
\begin{proof}
若 \(f(x)=-\infty\)，所需不等式没有内容。以下设 \(f(x)>-\infty\)。
先设 \(f\) 下半连续。给定 \(a<f(x)\)，开集 \(\{f>a\}\) 是 \(x\) 的邻域，所以存在
\(\alpha_0\) 使
\[
  \alpha\ge\alpha_0\quad\Longrightarrow\quad f(x_\alpha)>a.
\]
因此
\[
  \liminf_\alpha f(x_\alpha)
  =\sup_\beta\inf_{\alpha\ge\beta}f(x_\alpha)
  \ge \inf_{\alpha\ge\alpha_0}f(x_\alpha)\ge a.
\]
若 \(f(x)\in\R\)，令 \(a\uparrow f(x)\)；若 \(f(x)=+\infty\)，则上式对每个 \(a\in\R\)
成立。两种情形都给出所需不等式。

反过来，若 \(f\) 不下半连续，则存在 \(a\in\R\)，使 \(G=\{f>a\}\) 不是开集。取
\(x\in G\setminus\operatorname{int}G\)。每个 \(x\) 的邻域 \(U\) 都含有 \(y\) 满足 \(f(y)\le a\)。
令
\[
 J=\{(U,y):U\in\mathcal N(x),\ y\in U,\ f(y)\le a\},
\]
按邻域缩小排序，并令网值为 \(y\)。成对索引构造给出 \(y_j\to x\)。又因每个网值都满足
\(f(y_j)\le a\)，对每个阈值 \(j_0\) 都有
\[
  \inf_{j\ge j_0}f(y_j)\le a,
  \qquad
  \liminf_jf(y_j)=\sup_{j_0}\inf_{j\ge j_0}f(y_j)\le a<f(x),
\]
违背所假设的网不等式。
\end{proof}

\begin{remark}
若 \(X\) 第一可数，\cref{proposition:1-1-1-first-countable-sequential-tests} 中的邻域基构造把上一定理的
逆向反例网替换成序列，所以只检验序列即可。一般空间不能这样替换；本小节最后的习题给出反例。
\end{remark}

\begin{proposition}[上确界与有限和的稳定性]\label{proposition:1-2-3-2}
任意非空族 \(f_i:X\to\overline\R\) 的下半连续函数之逐点上确界仍下半连续。若
\(f_1,\ldots,f_n:X\to(-\infty,+\infty]\) 下半连续，则其逐点和也下半连续。
\end{proposition}
\begin{proof}
对上确界，
\[
 \left\{x:\sup_i f_i(x)>a\right\}
 =\bigcup_i\{x:f_i(x)>a\},
\]
右端开，故上确界下半连续。

只需证明两个函数之和。设 \(x_\alpha\to x\)。若 \(f(x),g(x)\) 都有限，给定 \(\varepsilon>0\)，
由下半连续性，
\[
 U_f=\{y:f(y)>f(x)-\varepsilon\},\qquad
 U_g=\{y:g(y)>g(x)-\varepsilon\}
\]
都是 \(x\) 的开邻域。收敛分别给出阈值 \(\alpha_f,\alpha_g\)，再由指标集的有向性取共同后继
\(\alpha_0\ge\alpha_f,\alpha_g\)。于是 \(\alpha\ge\alpha_0\) 时同时有
\[
 f(x_\alpha)>f(x)-\varepsilon,\qquad
 g(x_\alpha)>g(x)-\varepsilon.
\]
故 \(\liminf_\alpha(f+g)(x_\alpha)\ge f(x)+g(x)-2\varepsilon\)，再令
\(\varepsilon\downarrow0\)。若 \(f(x)=+\infty\)，则 \(g(x)>-\infty\)。给定 \(M\in\R\)，可取
实数 \(a,b\) 使 \(a+b=M\)、\(a<f(x)\) 且 \(b<g(x)\)：若 \(g(x)<+\infty\)，先取
\(b<g(x)\) 再令 \(a=M-b\)；若 \(g(x)=+\infty\)，可取 \(a=b=M/2\)。下半连续性给出共同尾部，
具体地，\(\{f>a\}\) 与 \(\{g>b\}\) 是 \(x\) 的开邻域；分别取收敛阈值并取二者的共同后继，
便使 \(f(x_\alpha)>a\) 且 \(g(x_\alpha)>b\)，从而和大于 \(M\)。因此和的 \(\liminf\) 为
\(+\infty\)。\(g(x)=+\infty\) 时同理。值域排除了 \(-\infty\)，所以没有
\(+\infty+(-\infty)\) 未定式。归纳得到有限和结论。
\end{proof}

\begin{theorem}[紧集上的极值]\label{theorem:1-2-3-3}
若 \(K\) 紧，\(f:K\to(-\infty,+\infty]\) 下半连续且不恒为 \(+\infty\)，则 \(f\) 在 \(K\) 上有限地
取到最小值。
\end{theorem}
\begin{proof}
令 \(m=\inf_Kf\)。先证 \(m>-\infty\)。若 \(m=-\infty\)，则闭集
\[
 F_n=\{x\in K:f(x)\le-n\}
\]
非空且递减，因而具有有限交性质。紧性给 \(x\in\bigcap_nF_n\)，这要求有限实数 \(f(x)\) 小于所有
\(-n\)，与值域矛盾。

因 \(f\) 不恒为 \(+\infty\)，\(m<+\infty\)。闭集
\[
 E_n=\{x\in K:f(x)\le m+1/n\}
\]
非空且递减。紧性给 \(x_*\in\bigcap_nE_n\)。于是 \(f(x_*)\le m+1/n\) 对所有 \(n\) 成立，故
\(f(x_*)\le m\)。下确界定义给相反不等式，因此 \(f(x_*)=m\in\R\)。
\end{proof}

\begin{proposition}[无选择的近似极小网]\label{proposition:1-2-3-minimizing-net}
设 \(F:X\to(-\infty,+\infty]\) 且 \(m=\inf_XF\in\R\)。存在一个网 \((x_j)\) 满足
\(F(x_j)\to m\)，其构造不需为每个误差预选一点。
\end{proposition}
\begin{proof}
令
\[
 J=\{(\varepsilon,x):\varepsilon>0,\ F(x)\le m+\varepsilon\},
\]
并规定
\[
 (\varepsilon,x)\preceq(\delta,y)\quad\Longleftrightarrow\quad\delta\le\varepsilon.
\]
下确界定义保证对每个 \(\varepsilon>0\) 都有可行点，所以 \(J\ne\varnothing\)。两个指标以
\(\eta=\min\{\varepsilon,\delta\}\) 加严，再取任一满足 \(F(z)\le m+\eta\) 的可行对，得到共同后继。
令网值 \(x_{(\varepsilon,x)}=x\)。给定 \(\rho>0\)，在任何误差坐标不超过 \(\rho\) 的指标之后，
\[
 m\le F(x_j)\le m+\rho.
\]
所以 \(F(x_j)\to m\)。指标集纳入全部可行对，没有指定整族 \(x_\varepsilon\)。
\end{proof}

\begin{theorem}[直接法原型]\label{theorem:1-2-3-4}
设 \(F:X\to(-\infty,+\infty]\) 下半连续，\(m=\inf_XF\in\R\)。若有网 \((x_\alpha)\) 满足
\(F(x_\alpha)\to m\)，且它的某个尾部落在一个相对紧集合中，则 \(F\) 取到 \(m\)。
\end{theorem}
\begin{proof}
设原网的指标集为 \(D\)。取阈值 \(\alpha_0\)，使尾网值落在相对紧集 \(A\) 中。闭包
\(K=\overline A\) 紧。把
\(D_0=\{\alpha\in D:\alpha\ge\alpha_0\}\) 作为尾指标集。它是有向集：\(D_0\) 中任意两个指标在原
指标集中有共同后继，而该后继仍大于 \(\alpha_0\)。包含映射 \(\iota:D_0\to D\) 最终共尾；事实上，
给定 \(\eta\in D\)，取 \(\eta\) 与 \(\alpha_0\) 的共同后继 \(\zeta\)，则
\(\alpha\ge\zeta\) 蕴含 \(\iota(\alpha)\ge\eta\)。因此尾网本身是原网的子网。

把该尾网看成 \(K\) 中的网；网紧性定理给出子网
\(x_{\iota(\phi(\beta))}\to x\in K\)，其中 \(\phi:B\to D_0\) 最终共尾。复合
\(\iota\circ\phi:B\to D\) 仍最终共尾：给定 \(\eta\in D\)，先取
\(\zeta\in D_0\) 使 \(\zeta\ge\eta\)，再在 \(\phi\) 的某个阈值后有
\(\phi(\beta)\ge\zeta\)。故这个收敛子网也是原网的子网，数值网
\(F(x_\alpha)\to m\) 沿它仍趋于 \(m\)。由下半连续性的网判别，
\[
 F(x)\le\liminf_\beta F(x_{\iota(\phi(\beta))})=m.
\]
另一方面 \(m\le F(x)\)，所以 \(F(x)=m\)。证明没有使用 Hausdorff 性；若空间非 Hausdorff，候选极限可能
不唯一，但任一由上述论证得到的极限仍是最小点。
\end{proof}

\begin{corollary}[次水平紧版直接法]\label{corollary:1-2-3-infcompact}
设 \(F:X\to(-\infty,+\infty]\) proper，即不恒为 \(+\infty\)，并且下半连续。若存在
\(c>\inf_XF\) 使 \(\{F\le c\}\) 相对紧，则 \(\inf_XF\in\R\) 且 \(F\) 取到它。特别地，proper、
下半连续且次水平紧的函数有最小点。
\end{corollary}
\begin{proof}
集合 \(E=\{F\le c\}\) 非空，并由下半连续性知其闭。相对紧性说明 \(\overline E\) 紧，而
\(E=\overline E\)，故 \(E\) 紧。紧集极值定理给出 \(x_*\in E\)，使 \(F(x_*)\) 是 \(F|_E\) 的有限
最小值。对 \(x\notin E\)，有 \(F(x)>c\ge F(x_*)\)，所以 \(x_*\) 也是全局最小点。这同时排除了
\(\inf_XF=-\infty\)。
\end{proof}

\begin{figure}[htbp]
\centering
\begin{tikzpicture}[node distance=1.2cm and 1.25cm,>=Stealth,every node/.style={align=center}]
 \node[book node] (m) {近似最小\\\(F(x_\alpha)\to m\)};
 \node[book node,below=of m] (s) {相对紧尾部};
 \node[book strong node,right=of m] (h) {同一网的\\两项信息};
 \node[book node,right=of h] (x) {收敛子网\\\(x_{\alpha_\beta}\to x\)};
 \node[book node,below=of x] (l) {下半连续\\\(F(x)\le m\)};
 \node[book strong node,left=of l] (e) {\(F(x)=m\)};
 \draw[book arrow] (m)--(h);
 \draw[book arrow] (s)--(h);
 \draw[book arrow] (h)--node[above]{紧性}(x);
 \draw[book arrow] (m)|-(l);
 \draw[book arrow] (x)--(l);
 \draw[book arrow] (l)--(e);
\end{tikzpicture}
\caption{紧性产生候选点，下半连续性识别候选点。}\label{fig:1-2-3-direct-method}
\end{figure}

下半连续本身既不保证下有界，也不保证取极值。任意族上半连续函数的上确界通常不再上半连续。等强制性只
控制候选点不逃逸，不自动给出变化泛函之间的极限不等式。第~3~章的 Fatou 引理和第~5~章的概率测度
开集不等式将在相应概念定义后给出同方向的 \(\liminf\) 机制；本章不使用这两个尚未陈述的结论。

\begin{exercise}
直接由定义证明：\(f\) 下半连续当且仅当每个次水平集闭。
\end{exercise}
\begin{exercise}
计算网 \(a_{(m,n)}=(-1)^m+1/n\)（\(\N^2\) 逐坐标排序）的 \(\liminf,\limsup\)，并求其聚点集。
\end{exercise}
\begin{exercise}
证明紧空间上的上半连续函数取到最大值。
\end{exercise}
\begin{exercise}
在通常的 ZFC 语境中，令 \(I=\mathcal P(\N)\)、\(X=\{0,1\}^{I}\)，并令
\[
 \Sigma=\{x\in X:\{i:x_i=1\}\text{ 至多可数}\}.
\]
证明 \(\Sigma\) 含所有有限支撑点，因而稠密但不闭；再证明 \(\Sigma\) 中任意逐坐标收敛序列的极限仍在
\(\Sigma\)。由此说明闭约束能量 \(I_\Sigma\) 满足所有序列的 \(\liminf\) 判别，却不下半连续。
\end{exercise}

下一节把紧性放入函数空间：点态候选只有在共同的局部控制下，才会成为连续并一致的函数极限。
